\documentclass[11pt,letterpaper]{article}
% KJO_preamble.tex
% Shared LaTeX preamble for all KJO documentation documents.
% \input this file immediately after \documentclass{}.

% --- Encoding and fonts -------------------------------------------------------
\usepackage[utf8]{inputenc}
\usepackage[T1]{fontenc}
\usepackage{lmodern}
\usepackage{microtype}

% --- Page geometry ------------------------------------------------------------
\usepackage[
  letterpaper,
  top=1in, bottom=1in,
  left=1.25in, right=1in
]{geometry}

% --- Core utilities -----------------------------------------------------------
\usepackage{amsmath}
\usepackage{graphicx}
\usepackage{xcolor}
\usepackage{parskip}        % space between paragraphs; no indent
\usepackage{enumitem}
\usepackage{caption}
\usepackage{subcaption}
\usepackage{float}

% --- Tables -------------------------------------------------------------------
\usepackage{booktabs}
\usepackage{longtable}
\usepackage{array}
\usepackage{multirow}

% --- Code listings ------------------------------------------------------------
\usepackage{listings}
\lstset{
  basicstyle=\ttfamily\small,
  keywordstyle=\bfseries\color{KJOblue},
  commentstyle=\itshape\color{gray},
  stringstyle=\color{KJOgreen!70!black},
  breaklines=true,
  frame=single,
  framesep=4pt,
  xleftmargin=6pt,
  xrightmargin=4pt,
  numbers=left,
  numberstyle=\tiny\color{gray},
  numbersep=8pt,
  language=C++,
}

% --- TikZ / PGF ---------------------------------------------------------------
\usepackage{tikz}
\usetikzlibrary{
  arrows.meta,
  shapes.geometric,
  shapes.misc,
  shapes.symbols,
  positioning,
  fit,
  calc,
  backgrounds,
  decorations.pathreplacing,
  decorations.markings,
  matrix,
}
\usepackage{pgfplots}
\pgfplotsset{compat=1.18}

% --- Headers / footers --------------------------------------------------------
\usepackage{fancyhdr}
\usepackage{lastpage}
\pagestyle{fancy}
\fancyhf{}
\renewcommand{\headrulewidth}{0.4pt}
\renewcommand{\footrulewidth}{0.4pt}
% Per-document: \lhead{}, \rhead{}, \cfoot{} are set in the main .tex file.

% --- Section formatting -------------------------------------------------------
\usepackage{titlesec}
\titleformat{\section}{\large\bfseries\color{KJOblue}}{{\thesection}}{0.8em}{}[\titlerule]
\titleformat{\subsection}{\normalsize\bfseries\color{KJOblue!80!black}}{{\thesubsection}}{0.8em}{}
\titleformat{\subsubsection}{\normalsize\bfseries}{{\thesubsubsection}}{0.8em}{}

% --- Hyperlinks ---------------------------------------------------------------
\usepackage[
  colorlinks=true,
  linkcolor=KJOblue,
  urlcolor=KJOblue!70!black,
  citecolor=KJOblue,
  pdfborder={0 0 0},
]{hyperref}

% --- Color palette -----------------------------------------------------------
\definecolor{KJOblue}   {RGB}{ 30,  90, 160}   % RCU
\definecolor{KJOred}    {RGB}{180,  30,  30}   % EMU
\definecolor{KJOgreen}  {RGB}{ 20, 120,  60}   % GSEMU
\definecolor{KJOorange} {RGB}{200, 110,  20}   % GSE / pressurization
\definecolor{KJOgray}   {RGB}{100, 100, 100}   % Buses / links
\definecolor{KJOpurple} {RGB}{110,  40, 140}   % Shared libraries
\definecolor{KJObg}     {RGB}{248, 248, 252}   % Light background for code

% --- Convenience macros -------------------------------------------------------

% Hardware component name (small caps)
\newcommand{\hw}[1]{\textsc{#1}}

% Software module name (monospace)
\newcommand{\sw}[1]{\texttt{#1}}

% Command name (monospace bold)
\newcommand{\cmd}[1]{\texttt{\textbf{#1}}}

% Pin / signal name
\newcommand{\sig}[1]{\texttt{#1}}

% Note / callout box
\usepackage{tcolorbox}
\tcbuselibrary{skins}
\newtcolorbox{notebox}{
  enhanced,
  colback=KJOblue!5,
  colframe=KJOblue!40,
  fonttitle=\bfseries,
  title=Note,
  left=4pt, right=4pt, top=2pt, bottom=2pt,
}
\newtcolorbox{warningbox}{
  enhanced,
  colback=KJOred!5,
  colframe=KJOred!40,
  fonttitle=\bfseries\color{KJOred},
  title=Warning,
  left=4pt, right=4pt, top=2pt, bottom=2pt,
}

% Title page macro: \KJOtitlepage{title}{subtitle}{version}{date}{description}
\newcommand{\KJOtitlepage}[5]{%
  \begin{titlepage}
    \centering
    \vspace*{2cm}
    {\Huge\bfseries\color{KJOblue} #1 \par}
    \vspace{0.5cm}
    {\Large\color{KJOgray} #2 \par}
    \vspace{1.5cm}
    \rule{0.5\textwidth}{0.4pt}\par
    \vspace{1cm}
    {\large Ken Overton \par}
    \vspace{0.3cm}
    {\normalsize \textit{KJO Liquid Rocket Motor Control System} \par}
    \vspace{1.5cm}
    \begin{tabular}{ll}
      \textbf{Version:} & #3 \\[4pt]
      \textbf{Date:}    & #4 \\
    \end{tabular}
    \vspace{1.5cm}\par
    \begin{minipage}{0.75\textwidth}
      \centering
      \normalsize #5
    \end{minipage}
    \vfill
    {\small\color{KJOgray} \textcopyright\ 2025--2026 Ken Overton.
     All rights reserved. \par}
  \end{titlepage}%
}

\newcommand{\docversion}{0.1}
\newcommand{\docdate}{September 2026}
\lhead{\textsc{KJO Valve Cycler}}
\rhead{\textsc{User Guide} --- v\docversion}
\cfoot{\thepage\ of \pageref{LastPage}}
\hypersetup{pdftitle={KJO Valve Cycler User Guide}, pdfauthor={Ken Overton}}

\begin{document}
% KJO_tikz_styles.tex
% Shared TikZ node and edge styles for all KJO block diagrams.
% \input this file inside the document body before any tikzpicture environments.

\tikzset{
  % ── Hardware block (peripheral chip / PCB component) ──────────────────────
  hwblock/.style={
    rectangle, draw=KJOgray, line width=1pt,
    fill=KJOgray!10, rounded corners=4pt,
    minimum width=2.6cm, minimum height=0.8cm,
    font=\small, align=center, inner sep=4pt,
  },
  % Sub-variant: MCU (larger, blue border)
  mcublock/.style={
    hwblock,
    draw=KJOblue, fill=KJOblue!8,
    minimum width=3.0cm, minimum height=1.0cm,
    font=\small\bfseries,
  },
  % Sub-variant: power / battery
  powerblock/.style={
    hwblock,
    draw=KJOorange, fill=KJOorange!8,
  },
  % ── Software module box ────────────────────────────────────────────────────
  swmodule/.style={
    rectangle, draw=KJOblue, line width=0.8pt,
    fill=KJOblue!6, rounded corners=3pt,
    minimum width=2.8cm, minimum height=0.7cm,
    font=\small, align=center, inner sep=3pt,
  },
  % Sub-variant: shared library
  sharedlib/.style={
    swmodule,
    draw=KJOpurple, fill=KJOpurple!6,
    dashed,
  },
  % ── Bus / link annotation ──────────────────────────────────────────────────
  busline/.style={
    draw=KJOgray, line width=1.5pt,
    -{Stealth[length=6pt,width=4pt]},
  },
  bidbusline/.style={
    draw=KJOgray, line width=1.5pt,
    {Stealth[length=6pt,width=4pt]}-{Stealth[length=6pt,width=4pt]},
  },
  radioline/.style={
    draw=KJOblue!60, line width=1.5pt, dashed,
    {Stealth[length=6pt,width=4pt]}-{Stealth[length=6pt,width=4pt]},
  },
  gpioline/.style={
    draw=KJOgreen!60!black, line width=1.2pt,
    -{Stealth[length=5pt,width=4pt]},
  },
  % ── Flowchart shapes ──────────────────────────────────────────────────────
  flowproc/.style={
    rectangle, draw=KJOblue!70, line width=0.8pt,
    fill=KJOblue!8, rounded corners=3pt,
    minimum width=3.2cm, minimum height=0.7cm,
    font=\small, align=center, inner sep=3pt,
  },
  flowdec/.style={
    diamond, draw=KJOorange, line width=0.8pt,
    fill=KJOorange!8,
    minimum width=2.2cm, minimum height=0.9cm,
    font=\small, align=center, inner sep=1pt,
  },
  flowterminal/.style={
    rounded rectangle, draw=KJOblue, line width=1pt,
    fill=KJOblue!12,
    minimum width=2.8cm, minimum height=0.7cm,
    font=\small\bfseries, align=center, inner sep=3pt,
  },
  flowarrow/.style={
    draw=KJOgray, line width=0.9pt,
    -{Stealth[length=5pt,width=4pt]},
  },
  % ── Propellant schematic styles ────────────────────────────────────────────
  % Generic pipe line
  pipe/.style={
    draw=black, line width=1.8pt,
  },
  pipearrow/.style={
    pipe,
    -{Stealth[length=7pt,width=5pt]},
  },
  % Oxidizer (N2O) — blue
  oxpipe/.style={
    draw=KJOblue, line width=1.8pt,
  },
  oxpipearrow/.style={
    oxpipe,
    -{Stealth[length=7pt,width=5pt]},
  },
  % Fuel — red
  fuelpipe/.style={
    draw=KJOred, line width=1.8pt,
  },
  fuelpipearrow/.style={
    fuelpipe,
    -{Stealth[length=7pt,width=5pt]},
  },
  % GSE fill line — green
  gsepipe/.style={
    draw=KJOgreen!70!black, line width=1.8pt,
  },
  gsepipearrow/.style={
    gsepipe,
    -{Stealth[length=7pt,width=5pt]},
  },
  % Pressurization — orange dashed
  presspipe/.style={
    draw=KJOorange, line width=1.4pt, dashed,
  },
  presspipearrow/.style={
    presspipe,
    -{Stealth[length=6pt,width=4pt]},
  },
  % Tank (tall rectangle, rounded corners)
  tank/.style={
    rectangle, draw=black, line width=1.5pt,
    rounded corners=4pt, fill=white,
    minimum width=1.8cm, minimum height=2.4cm,
    font=\small\bfseries, align=center, inner sep=4pt,
  },
  n2otank/.style={
    tank, draw=KJOblue, fill=KJOblue!8,
  },
  fueltank/.style={
    tank, draw=KJOred, fill=KJOred!8,
  },
  srctank/.style={
    tank, draw=KJOgreen!70!black, fill=KJOgreen!6,
  },
  % Valve box (servo-actuated)
  valvebox/.style={
    rectangle, draw=black, line width=1pt,
    fill=white, minimum width=2.2cm, minimum height=0.65cm,
    font=\footnotesize\bfseries, align=center, inner sep=3pt,
  },
  gseval/.style={
    valvebox, draw=KJOgreen!70!black, fill=KJOgreen!8,
  },
  emuval/.style={
    valvebox, draw=KJOred, fill=KJOred!6,
  },
  % Component (generic inline component)
  component/.style={
    rectangle, draw=KJOgray, line width=0.8pt,
    fill=white!96!gray, minimum width=2.2cm, minimum height=0.6cm,
    font=\footnotesize, align=center, inner sep=3pt,
  },
  % Instrument callout circle (PT, TT, TC)
  instrument/.style={
    circle, draw=KJOgray, line width=0.8pt,
    fill=white, minimum size=0.6cm,
    font=\tiny\bfseries, align=center, inner sep=0pt,
  },
  % Sensor lead line
  sensorlead/.style={
    draw=KJOgray, line width=0.7pt,
  },
  % Boundary box (dashed outline for GSE / Rocket regions)
  gseboundary/.style={
    draw=KJOgreen!50, line width=1pt, dashed,
    rounded corners=6pt,
  },
  rocketboundary/.style={
    draw=KJOred!40, line width=1pt,
    rounded corners=6pt,
  },
  % Sequence diagram lifeline
  lifeline/.style={
    draw=KJOgray, line width=0.6pt, dashed,
  },
  % Sequence diagram message arrow
  seqarrow/.style={
    draw=KJOblue, line width=0.9pt,
    -{Stealth[length=5pt,width=4pt]},
  },
  % State node (FSM)
  statenode/.style={
    circle, draw=KJOblue, line width=1pt,
    fill=KJOblue!10,
    minimum size=1.2cm,
    font=\small\bfseries, align=center,
  },
  statetransition/.style={
    draw=KJOgray, line width=0.9pt,
    -{Stealth[length=5pt,width=4pt]},
    font=\footnotesize,
  },
}


\KJOtitlepage%
  {Valve Cycler}%
  {Bench Tool User Guide}%
  {\docversion}{\docdate}%
  {A standalone bench tool that repeatedly cycles a single servo-actuated
   valve between two PWM endpoints, for mechanical endurance/wear testing
   and for exercising a valve during bench assembly work.}

\tableofcontents\newpage

%═══════════════════════════════════════════════════════════════════════════════
\section{Purpose}
%═══════════════════════════════════════════════════════════════════════════════

Valve Cycler drives one servo-actuated valve back and forth between two
operator-supplied PWM counts (``Upper'' and ``Lower''), with a configurable
pause between moves, either a fixed number of times or indefinitely. It is
used for mechanical endurance testing, seating a newly-assembled valve, or
simply exercising a valve on the bench without needing the full production
firmware running.

%═══════════════════════════════════════════════════════════════════════════════
\section{Hardware Required}
%═══════════════════════════════════════════════════════════════════════════════

\begin{itemize}[nosep]
  \item Adafruit Feather M0 (ATSAMD21G18A)
  \item Adafruit 8-channel PWM/Servo FeatherWing (PCA9685), I\textsuperscript{2}C
  \item One servo-actuated valve, servo wired to \textbf{PCA9685 channel 7}
  \item USB cable for Serial (115200 baud)
\end{itemize}

\begin{notebox}
The servo channel (7) is hardcoded in the sketch. To cycle a valve wired to
a different channel, edit \texttt{SERVO\_CHANNEL} in \texttt{main.cpp} and
reflash.
\end{notebox}

%═══════════════════════════════════════════════════════════════════════════════
\section{How to Use}
%═══════════════════════════════════════════════════════════════════════════════

Open a Serial monitor at 115200 baud after powering the board. On every
power-up (and again after a run completes) the tool asks four questions,
each with a \texttt{y/n} confirmation before moving to the next:

\begin{enumerate}[nosep]
  \item \textbf{Upper} -- the maximum PWM count (one travel endpoint).
  \item \textbf{Lower} -- the minimum PWM count (the other travel endpoint).
  \item \textbf{Delay D} -- seconds to hold at each endpoint before moving
        to the other.
  \item \textbf{Number of cycles} -- how many Lower$\to$Upper$\to$Lower
        round trips to run, or \textbf{$-1$} for indefinite cycling.
\end{enumerate}

If you answer \texttt{n} to any confirmation, the tool re-asks that same
question. Once all four are confirmed, cycling starts immediately: the
servo moves to Lower, waits \texttt{D} seconds, moves to Upper, waits
\texttt{D} seconds, and repeats. A finite run prints \texttt{"N cycles
completed."} and returns to the setup questions for another run; an
indefinite run (\texttt{$-1$}) never stops on its own -- power-cycle the
board to end it.

%═══════════════════════════════════════════════════════════════════════════════
\section{Notes}
%═══════════════════════════════════════════════════════════════════════════════

\begin{itemize}[nosep]
  \item The Serial console has no timeout -- the tool waits indefinitely
        at every prompt.
  \item Backspace is supported while typing a value.
  \item Upper and Lower are used exactly as entered -- the tool does not
        check that Upper is actually greater than Lower, or that either
        value is within the servo's safe travel range. Use the companion
        \textbf{Servo Tuner} bench tool first to find safe endpoints for a
        given valve.
\end{itemize}

%═══════════════════════════════════════════════════════════════════════════════
\section{Version History}
%═══════════════════════════════════════════════════════════════════════════════
\begin{center}
\begin{tabularx}{\textwidth}{@{}llY@{}}
\toprule
Version & Date & Notes \\
\midrule
0.1 & September 2026 & Initial user guide. \\
\bottomrule
\end{tabularx}
\end{center}

\end{document}
